\section{Corporate issuers, ownership and governance}

\begin{qfdefinition}[Organisational forms and corporate issuers: sole proprietorships, partnerships, limited-liability hybrids and corporations compared on liability, taxation, life, transferability and access to capital, and public against private issuers]\label{def:organisational-forms-issuers}
\qffield{Statement} A business can be organised as a sole proprietorship, a partnership, a corporation or a hybrid such as a limited liability company (LLC) or limited liability partnership (LLP) \citep[sec.~2.1]{principles-of-finance-2026}. The forms differ in five ways:
\begin{itemize}
\item \emph{Liability.} A proprietor, or a partner in a general partnership, is personally liable for all of the business's debts. A corporation is a legal entity separate from its owners, so shareholders can lose only what they invested. The three limited-liability partnership-type forms differ; Delaware law illustrates them.
\begin{itemize}
\item A \emph{limited partnership} (LP) has at least one general partner, liable as a partner in a general partnership \citep[sec.~17-403(b)]{delaware-partnership-llc-2026}, and limited partners, who are not liable for its obligations unless they are also general partners or take part in control of the business; even then they are liable only to those who reasonably believed them to be general partners (sec.~17-303(a)).
\item A \emph{limited liability partnership} (LLP) is a general partnership that has filed a statement of qualification (sec.~15-1001). It has no limited partners: obligations arising while it is an LLP, ``whether arising in contract, tort or otherwise, [are] solely the obligation of the partnership'', and a partner is not personally liable for them ``solely by reason of being or so acting as a partner'' (sec.~15-306(c)). The shield is qualified: partners can agree to be liable (sec.~15-306(e)), and it reaches only obligations arising or incurred while the partnership is an LLP. Because it removes only liability that arises by reason of being a partner, a partner who personally commits a wrong or gives a guarantee is, on that reading (Interpretation), liable on that separate ground.
\item A \emph{limited liability company} (LLC) gives the same form of shield to its members and managers, who may also agree to be personally obligated (sec.~18-303).
\end{itemize}
\item \emph{Taxation.} Proprietorship and partnership profits are taxed once, as the owners' income. Corporate profits are taxed at the company and again when paid out as dividends, which is double taxation.
\item \emph{Life and transfer.} A corporation's life is unlimited and its shares transfer easily; a proprietorship ends with its owner, and a partnership interest is hard to sell.
\item \emph{Capital.} Unlimited liability and illiquid interests make large amounts of capital hard to raise. The corporation's advantages are why, by value of sales, more than 80 percent of US business is done by corporations (same section).
\end{itemize}

In the United States an S corporation is taxed like a partnership and is exempt from corporate income tax, but may have no more than 100 shareholders. LLCs give owners liability protection with partnership-style taxation (same section). The Internal Revenue Service describes the LLC as a business structure allowed by state statute and the S corporation as an election to pass income, losses, deductions and credits through to shareholders \citep[``Business Structures''; ``S Corporations'']{irs-business-structures-2026}.

\emph{Features of corporate issuers.} Equity is the residual claim: it has a significant role in management and is perpetual. Debt is a fixed claim with little role in management and tax-deductible interest \citep[pp.~1--2]{applied-corporate-finance-2014}. Ownership is separated from control, with shareholders electing a board that hires managers. The resulting agency problems are the subject of \hyperref[meth:corporate-governance-evaluation]{governance}.

\emph{Public against private issuers.} A public company's shares are listed and traded. It gains access to a deep pool of capital, a market price and liquidity for its owners. In exchange it bears listing and disclosure costs, investor relations, and possible loss of control to dispersed holders and bidders \citep[p.~302]{applied-corporate-finance-2014}. A US issuer must register a class of equity under the Exchange Act within 120 days after the end of a fiscal year in which it has total assets above \$10 million and the class is held of record by 2,000 persons, or by 500 who are not accredited investors \citep[subsec.~(g)(1)(A)]{usc-15-78l}. Registration brings periodic reporting even if the shares are not listed.

A private company raises equity from founders, venture capital and private equity, and debt from banks or private placements. Its shares cannot be resold easily, its owners are often undiversified, and its disclosure is limited to what lenders and investors negotiate.

\qffield{Conditions} The categories and tax treatments are US law as described by OpenStax and the IRS in 2026. The liability rules quoted are Delaware's, read in the official online code in October 2026; other states and countries differ, for example in whether an LLP shield covers contractual as well as tort claims and in the limited partner's control rule. Other jurisdictions have parallel forms, such as the private limited company and the public limited company, with different rules on minimum capital, disclosure and the taxation of distributions. Partial relief from double taxation, through imputation credits or lower dividend rates, is described under \hyperref[meth:dividend-policy-evaluation]{dividend policy}.

\qffield{Practice} \emph{Original.} A business earns 100 before tax. As a corporation with a 25~percent corporate rate whose owner pays 30~percent on dividends, a full payout leaves the owner $100\times0.75\times0.70=52.50$, an effective rate of 47.5~percent. As a pass-through taxed at the owner's 40~percent rate it leaves 60.

If the corporation retains the profit, 75 stays in the business and personal tax is deferred until distribution or sale. The corporate form then suits a business that reinvests and needs outside equity. The pass-through form suits one that distributes most of its profit to a few owners. The comparison changes with each jurisdiction's rates and with whether personal liability matters.

\qffield{Pitfalls} Treating limited liability as costless: lenders to small companies often require personal guarantees, which restore the owner's liability. Comparing forms on tax alone, ignoring the cost of raising capital and the value of transferable shares. Equating ``private'' with ``small'': large private-equity-owned groups are private issuers, and some public companies are small. Assuming that an unlisted company has no reporting duty; holder-of-record thresholds can require registration.

\qffield{See also} \hyperref[def:stakeholder-claims]{stakeholder claims}; \hyperref[def:equity-instruments]{equity instruments} (listed against private equity); \hyperref[meth:cost-of-debt-equity-estimation]{private-company cost of capital}.

\end{qfdefinition}

\begin{qfdefinition}[Claims and interests of lenders, shareholders and other stakeholders: fixed against residual claims, risk shifting under limited liability, conflicts with managers, employees, customers, suppliers, governments and society, and the ESG factors that describe an issuer]\label{def:stakeholder-claims}
\qffield{Statement} \emph{Lenders and shareholders.} Lenders hold a fixed claim: contractual interest and principal, paid before shareholders, with covenants and collateral to protect it. Shareholders hold the residual claim: whatever is left after all fixed claims. With limited liability, that claim is like a call option on the firm's assets \citep[pp.~1--2, 26--27]{applied-corporate-finance-2014}. Their payoffs therefore differ in shape. Lenders gain nothing from upside beyond full repayment but bear losses in default. Shareholders keep all the upside and lose at most their investment.

Three actions transfer wealth from lenders to shareholders after the debt is priced: taking riskier projects, adding debt of equal or higher priority, and paying large dividends or buying back shares (pp.~26--27). Bond prices fall when such actions are announced. Damodaran cites a fall of about 19 percent in Nabisco bond prices on the leveraged buyout (p.~26). Lenders respond with covenants restricting investment, dividends and further borrowing, with security, convertibles or puttable bonds, and with higher interest rates (pp.~41--42).

\emph{Other stakeholders.} Managers, employees, customers, suppliers, governments and society also have interests in the firm (pp.~11--12):
\begin{itemize}
\item managers want pay, security and control;
\item employees want wages, job security and safe conditions;
\item customers want quality, price and continued service;
\item suppliers want payment and continued business;
\item governments want taxes and compliance;
\item society wants the firm to bear the costs it imposes, such as pollution.
\end{itemize}
The G20/OECD Principles ask that stakeholders' rights established by law or mutual agreement be respected. Bondholders' rights should be protected, including through a trustee \citep[principles~VI.D and~VI.D.6, pp.~48--49]{oecd-g20-principles-2023}.

\emph{ESG factors of issuers.} Environmental factors are an issuer's impacts on, and dependencies on, nature: emissions, energy and water use, waste and pollution, and exposure to climate hazards. Social factors cover its workforce, its supply chain, its customers and product safety, and its communities. Governance factors cover its board, its ownership and control, its pay, its accounting and its controls. The EU standards define sustainability matters through \emph{double materiality}. A matter is material if the company's impacts on people or the environment are material, or if the matter is expected to have material financial effects on the company, or both \citep[Annex, ESRS~1 paras.~22--25]{eu-2026-1563-esrs}.

\qffield{Conditions} The claims described are those of an ordinary going concern. In bankruptcy the order of claims is set by insolvency law and contracts, and shareholders usually receive nothing until creditors are paid. Which stakeholder interests directors must consider depends on company law: some jurisdictions oblige directors to have regard to employees and the community, while others frame the duty around shareholder value \citep[pp.~33--35]{applied-corporate-finance-2014}.

\qffield{Practice} \emph{Original: risk shifting.} A firm owes 100 due in one year. A safe project leaves it worth 60 or 140 with equal probability; a risky one leaves it worth 20 or 180 (expected values are compared undiscounted for clarity).
\begin{itemize}
\item Safe: equity is worth $\frac12(0)+\frac12(40)=20$ and debt $\frac12(60)+\frac12(100)=80$.
\item Risky, same mean 100: equity becomes $\frac12(0)+\frac12(80)=40$ and debt $\frac12(20)+\frac12(100)=60$. Twenty moves from lenders to shareholders with no change in firm value.
\item A risky project worth 10 or 180 (mean 95) still gives equity 40. Shareholders would choose it although it destroys 5 of value, because lenders bear the loss of 25.
\end{itemize}
Lenders anticipating this charge more or write covenants, so the cost comes back to shareholders as a higher cost of debt or forgone flexibility.

\qffield{Pitfalls} Treating ``maximise shareholder value'' as costless to other claimants: it is a sensible objective only when lenders are protected, information is honest and social costs are borne by the firm \citep[pp.~13--14]{applied-corporate-finance-2014}. Adding stakeholder objectives without a way to weigh them leaves managers accountable to no one (p.~33). Counting an ESG factor as material because it is popular rather than because it affects people, the environment or cash flows.

\qffield{See also} \hyperref[def:organisational-forms-issuers]{organisational forms}; \hyperref[meth:corporate-governance-evaluation]{governance}; \hyperref[meth:esg-exposure-evaluation]{ESG exposures}; \hyperref[mod:capital-structure-theory]{agency costs of debt}.

\end{qfdefinition}

\begin{qfmethod}[Corporate governance: principal--agent conflicts, internal and external governance mechanisms, the risks of weak and benefits of effective governance, and evaluating a company's governance]\label{meth:corporate-governance-evaluation}
\qffield{Statement} \emph{Conflicts.} Separation of ownership from control creates principal--agent conflicts between shareholders and managers. Managers may pursue pay, perquisites, empire-building acquisitions and entrenchment, or may hold back bad news \citep[pp.~14--19, 27--29]{applied-corporate-finance-2014}. Further conflicts arise:
\begin{itemize}
\item between controlling and minority shareholders, over private benefits, related-party transactions and issuance on favourable terms \citep[principle~II.G, pp.~20--21]{oecd-g20-principles-2023};
\item between shareholders and lenders (\hyperref[def:stakeholder-claims]{stakeholder claims});
\item between the firm and society, over costs it does not bear.
\end{itemize}

\emph{Mechanisms.} Internal mechanisms include:
\begin{itemize}
\item shareholder voting on directors, pay and major transactions;
\item a board with independent members and audit, remuneration and nomination committees;
\item executive pay linked to long-term value;
\item internal controls and audit.
\end{itemize}
The Principles recommend separating the chair and chief executive, or appointing a lead independent director, and committees with independent members \citep[principles~V.E and~V.E.2, pp.~39--41]{oecd-g20-principles-2023}. Related-party transactions should be approved with the abstention of interested parties \citep[principle~II.F, pp.~19--20]{oecd-g20-principles-2023}.

External mechanisms include the market for corporate control, activist investors, lenders' covenants, auditors, regulators, disclosure rules and the courts \citep[pp.~38--42]{applied-corporate-finance-2014}. In the EU, institutional investors and asset managers must publish an engagement policy or explain why not. They must also disclose their voting, including significant votes and their use of proxy advisers. Shareholders vote on the remuneration policy \citep[Arts.~3g and~9a of Directive 2007/36/EC as inserted]{eu-2017-828-srd2}.

\emph{Risks and benefits.} Weak governance leads to value-destroying acquisitions, excessive pay, misreporting, expropriation of minorities and a higher cost of capital. Effective governance aligns decisions with long-term value, but its payoff is hard to measure. Gompers, Ishii and Metrick's index of 24 provisions produced a long-short return of 8.5 percent a year in the 1990s. Later studies tie governance more to valuation levels than to operating performance, and find little evidence that board composition alone creates value \citep[pp.~25--26]{applied-corporate-finance-2014}. Takeover threats discipline managers: targets tend to have underperformed, and state laws restricting takeovers lowered stock prices (pp.~39--40).

\qffield{Notation} A governance evaluation asks four questions (Interpretation):
\begin{enumerate}
\item Who controls the board and the votes, and how far does control exceed economic ownership?
\item Can shareholders replace directors and managers, and how costly is a change of control?
\item Are managers paid for long-term value, and is pay disclosed and voted on?
\item Are reporting, audit and related-party transactions independent of the people they constrain?
\end{enumerate}
For an investor the effectiveness question is whether a weakness changes expected cash flows or the risk that value is expropriated.

\qffield{Practice} \emph{Original.} A listed company has a founder chief executive who also chairs a nine-member board. Three directors are independent, and the audit committee includes one executive. The founder holds 15 percent of the equity but 64 percent of the votes through a class with ten votes a share (\hyperref[def:global-ownership-structures]{ownership structures}). The company has bought three businesses from the founder's family in five years, and pay is mostly an annual cash bonus on revenue growth.

Evaluation:
\begin{itemize}
\item \emph{Control.} Votes exceed ownership by more than four times, so outside shareholders cannot replace the board, and the market for control is shut.
\item \emph{Related parties.} The purchases from the family test whether independent directors approved them and on what valuation. Without abstention and fairness evidence, assume a transfer risk.
\item \emph{Pay.} Revenue-based bonuses reward growth regardless of return on capital, and the acquisitions fit that incentive.
\item \emph{Board.} An executive on the audit committee and a combined chair and chief executive weaken monitoring.
\end{itemize}
In valuation the response is to haircut cash flows for likely value-destroying acquisitions, to apply a lower multiple than otherwise comparable peers, or both. Where the founder's long-term stake aligns incentives, record the offsetting benefit rather than scoring every concentration as negative.

\qffield{Conditions} Governance codes are mostly comply-or-explain, and the same practice can be appropriate in one company and harmful in another. The G20/OECD Principles are a non-binding international standard \citep[pp.~15--21, 39--41]{oecd-g20-principles-2023}. Studies of governance and returns are period- and market-specific, and governance scores combine judgements that may not match an investor's view.

\qffield{Pitfalls} Scoring governance by a checklist of board features alone. Treating a controlling shareholder as bad in itself: concentration can solve the monitoring problem it also creates \citep[principle~II.G, p.~20]{oecd-g20-principles-2023}. Counting a governance weakness in both the discount rate and the cash flows. Assuming that activist pressure always improves value; it can also be short-term or conflicted, for example through empty voting \citep[pp.~42--43]{applied-corporate-finance-2014}.

\qffield{See also} \hyperref[def:global-ownership-structures]{ownership structures}; \hyperref[def:equity-rights-voting]{shareholder rights and voting}; \hyperref[meth:esg-exposure-evaluation]{ESG exposures}; \hyperref[meth:reporting-warning-signs]{reporting warning signs}.

\end{qfmethod}

\begin{qfdefinition}[Ownership structures around the world: dispersed and concentrated ownership, institutional, public-sector, corporate and family owners, dual-class shares, pyramids and cross-holdings, and their governance consequences]\label{def:global-ownership-structures}
\qffield{Statement} Ownership of listed companies falls into two broad patterns. \emph{Dispersed} ownership, where no holder controls the company, makes the conflict between managers and shareholders central. \emph{Concentrated} ownership, where a family, a corporation or the state controls, reduces that conflict through monitoring. It creates another, between the controller and minority holders \citep[principle~II.G, p.~20]{oecd-g20-principles-2023}.

At end-2017 the 10,000 largest listed companies, about 90 percent of world market value, were owned as follows \citep[pp.~5--6]{oecd-owners-listed-2019}:
\begin{itemize}
\item institutional investors held 41 percent of global market capitalisation;
\item the public sector held 14 percent, about \$10 trillion;
\item private corporations, including holding companies, together with strategic individuals and families, held 18 percent;
\item in half of the companies, the three largest holders owned more than 50 percent, and in three quarters more than 30 percent;
\item private corporations held 40 to 55 percent of listed equity in Chile, the Philippines and Turkey.
\end{itemize}
Control can exceed cash-flow ownership through four devices \citep[pp.~16--17, 23--24]{applied-corporate-finance-2014}:
\begin{itemize}
\item shares with different voting rights (dual-class or non-voting preferred);
\item pyramids, in which control passes through a chain of holding companies;
\item cross-holdings among group companies, such as the keiretsu or Tata's 31.83 percent cross-holdings;
\item golden shares that give a government a veto.
\end{itemize}
Damodaran describes market-based systems, bank- and group-centred systems (German and Japanese) and state-led systems as alternative governance models (pp.~34--35).

\qffield{Notation} With $n_A$ shares of one vote and $n_B$ shares of $v$ votes, a holder of all class~B has vote share $vn_B/(vn_B+n_A)$ and cash-flow share $n_B/(n_A+n_B)$. In a pyramid, cash-flow rights multiply along the chain: owning a fraction $a$ of a holding company that owns $b$ of the operating company gives cash-flow rights $ab$, while control follows the weakest link.

\qffield{Practice} \emph{Original.}
\begin{itemize}
\item \emph{Dual class.} A founder holds 15 million class~B shares with ten votes each; the public holds 85 million class~A shares with one vote. The founder has 15 percent of the cash flows and $150/(150+85)=63.83$ percent of the votes.
\item \emph{Pyramid.} A family owns 51 percent of a holding company that owns 51 percent of an operating company. It controls the operating company with $0.51\times0.51=26.01$ percent of its cash flows. If the operating company sells an asset to a family-owned business at 10 below value, the family bears 2.60 of the loss through its stake and gains 10.
\end{itemize}
The wedge between the two is what makes tunnelling profitable, and what minority protections are meant to check. Those protections include independent approval of related-party transactions, pre-emptive rights and qualified majorities \citep[principles~II.F--II.G, pp.~19--21]{oecd-g20-principles-2023}.

\qffield{Conditions} The ownership figures are for end-2017 and classify holders by the type of the ultimate owner; they say nothing about each holder's voting behaviour. Many institutional investors are index funds whose engagement is through voting rather than trading. Index providers exclude strategic and restricted holdings from free float, which matters for liquidity (\hyperref[def:equity-issuance-trading-liquidity]{free float}).

\qffield{Pitfalls} Reading a large institutional share as a sign of active monitoring. Ignoring control-enhancing devices when calculating a controller's incentive. Assuming that a state-controlled company maximises shareholder value: the state may pursue policy goals through it.

\qffield{See also} \hyperref[meth:corporate-governance-evaluation]{governance}; \hyperref[def:equity-rights-voting]{voting rights}; \hyperref[meth:valuation-approach-choice]{conglomerate discounts}.

\end{qfdefinition}

\begin{qfmethod}[Identifying and evaluating ESG exposures and opportunities: double materiality, impacts, risks and opportunities, divergent ratings, and translating material factors into cash flows, risk and value]\label{meth:esg-exposure-evaluation}
\qffield{Statement} \emph{Identifying exposures.} Under the EU sustainability standards a company assesses sustainability matters for double materiality:
\begin{itemize}
\item \emph{Impact materiality}: actual or potential, positive or negative impacts on people or the environment, in its own operations and value chain, over the short, medium or long term \citep[Annex, ESRS~1 paras.~38--44]{eu-2026-1563-esrs}. Severity combines scale, scope and irremediability; potential impacts weigh severity against likelihood.
\item \emph{Financial materiality}: matters that trigger or may trigger material financial effects on the company's development, financial position, performance, cash flows, access to finance or cost of capital (paras.~45--49). Sources include impacts that rebound on the company, dependencies on natural, human and social resources, and other factors such as climate hazards and regulatory change.
\end{itemize}
The assessment starts top-down from the company's strategy, business model, sector, geography and value chain (para.~27). An investor uses the same logic but asks only about financial effects, including those that come back from impacts on others. The G20/OECD Principles treat sustainability information as material if a reasonable investor would consider it important for investment or voting decisions; a company's own impacts can be material, for example where emissions are capped or taxed \citep[principle~VI.A.1, p.~45]{oecd-g20-principles-2023}.

\emph{Evaluating exposures.} Translate each material factor into the inputs of a valuation \citep[Interpretation]{applied-corporate-finance-2014}:
\begin{itemize}
\item revenue, for products whose demand shifts with regulation or preference;
\item costs, such as carbon prices, remediation and labour standards;
\item capital expenditure, for adaptation or transition;
\item contingent liabilities, for litigation and recalls;
\item the probability of large losses;
\item the cost of capital, but only for risk that is systematic and not already in the cash flows.
\end{itemize}
Damodaran finds no consistent return difference for socially conscious firms but notes that markets punish firms that impose large social costs \citep[pp.~44--45]{applied-corporate-finance-2014}.

\emph{Ratings.} Third-party ESG ratings disagree. Across six raters, correlations ran from 0.38 to 0.71. Of the divergence, 56 percent came from measuring the same attribute differently, 38 percent from scope (which attributes are included) and 6 percent from weights \citep[abstract and sec.~1]{aggregate-confusion-2022}. A rater's overall view also colours how it scores individual categories (abstract). The authors advise investors to use indicator-level data or a rater whose definitions match their objective (sec.~7).

\qffield{Practice} \emph{Original, Harbour Tools} (\hyperref[ex:harbour-tools-statements]{case}). Material matters for a power-tool maker include:
\begin{itemize}
\item product safety and recalls (social, cash flows);
\item labour standards at contract manufacturers (social, supply risk);
\item energy use and emissions in its plants (environmental, costs);
\item battery collection and recycling obligations (environmental, regulation);
\item the governance of its acquisitions (governance, capital allocation).
\end{itemize}

One factor quantified, with illustrative inputs. Scope~1 and~2 emissions are 30,000 tonnes a year. Carbon costs reach Harbour through energy prices at EUR~130 a tonne from 2030, and 40 percent is passed on to customers. The pre-tax cost is $30{,}000\times130\times0.6=$ EUR~2.34 million a year. After tax at 25 percent, capitalised at 8~percent less 2.5~percent growth and discounted four years to the end of 2025, it is worth about 23.45 million, or 0.23 a share. That is 2.7 percent of the 8.49 value per share (\hyperref[meth:forecast-horizon-steady-state]{Harbour valuation}).

The exposure is real but modest, and it belongs in the cash flows. Adding a premium to the discount rate as well would count it twice. Opportunities, such as demand for efficient cordless tools, enter revenue forecasts in the same way.

\qffield{Conditions} ESRS~1 as amended in 2026 applies to financial years beginning on or after 1~January 2027, with optional earlier use \citep[Arts.~2--3]{eu-2026-1563-esrs}. Other regimes, such as the ISSB standards, use financial materiality only. The carbon example's prices and pass-through are assumptions, not forecasts.

\qffield{Pitfalls} Relying on one ESG rating as a measure of exposure when raters disagree on what they measure. Treating every sustainability matter as financially material, or none. Adjusting the discount rate for diversifiable or already-modelled risks. Ignoring positive impacts and opportunities, which ESRS does not net against negative impacts.

\qffield{See also} \hyperref[def:stakeholder-claims]{stakeholders and ESG factors}; \hyperref[meth:corporate-governance-evaluation]{governance}; \hyperref[meth:company-position-analysis]{company analysis}.

\end{qfmethod}
